\section{The 1:10 diluted-target experiment}
\label{app:dilute10}

This appendix gives the values behind Section~\ref{sec:results-dilute} at four decimals, and the
run facts the text cites. Every value comes from the experiment's analysis file or from the run
records named in Appendix~\ref{app:repro}; the methods are in Section~\ref{sec:dilute-methods}
and the pre-registration in Appendix~\ref{sec:prereg-dilute}.

\paragraph{Per-fold values and the ten intervals.} Table~\ref{tab:dl-perfold} gives every arm's
diluted values per fold for the four metrics of Appendix~\ref{app:perfold}, and
Table~\ref{tab:dl-contrasts4} the \dlNtWord{} paired $t$-intervals: the \dlNtG{} forms of $G$ and
the \dlNtD{} of $D$. The matched-step forms of $G$ give the same branches as the primary forms.

% GENERATED FILE -- do not edit by hand.
% Rebuild:  python make_all.py (generator: make_dilute10_tables.py)
% Appendix K: per-fold diluted values. runs/dilute10_analysis/results.json per_arm.
\begin{table}[htbp]
\centering
\caption{\textbf{Per-fold values of the diluted runs}, cleaned pool, at each arm's read-off step, with the fold-mean and sample sd. Completeness at 24\% purity and purity at 50\% completeness use the prevalence correction of Equation~\eqref{eq:purity}.}
\label{tab:dl-perfold}
\footnotesize
\setlength{\tabcolsep}{4pt}
\begin{tabular}{llrrrrrrr}
\toprule
Metric & Arm & Step & fold 1 & fold 2 & fold 3 & fold 4 & mean & sd \\
\midrule
TPR@$10^{-3}$ & DANN & 4000 & $0.5306$ & $0.5204$ & $0.5102$ & $0.5155$ & $0.5192$ & $0.0087$ \\
 & BNM & 4000 & $0.3980$ & $0.5408$ & $0.2857$ & $0.2887$ & $0.3783$ & $0.1203$ \\
 & CORAL & 2000 & $0.2449$ & $0.2041$ & $0.1531$ & $0.1753$ & $0.1943$ & $0.0397$ \\
 & FreeMatch & 1500 & $0.2959$ & $0.6327$ & $0.3469$ & $0.4742$ & $0.4374$ & $0.1502$ \\
\midrule
AUROC & DANN & 4000 & $0.9510$ & $0.8971$ & $0.8808$ & $0.7972$ & $0.8815$ & $0.0637$ \\
 & BNM & 4000 & $0.9445$ & $0.9804$ & $0.9218$ & $0.9096$ & $0.9391$ & $0.0311$ \\
 & CORAL & 2000 & $0.7813$ & $0.8089$ & $0.7990$ & $0.7610$ & $0.7876$ & $0.0211$ \\
 & FreeMatch & 1500 & $0.9660$ & $0.9874$ & $0.9731$ & $0.9781$ & $0.9762$ & $0.0090$ \\
\midrule
completeness at 24\% purity & DANN & 4000 & $0.5102$ & $0.5204$ & $0.4490$ & $0.4639$ & $0.4859$ & $0.0348$ \\
 & BNM & 4000 & $0.3163$ & $0.5204$ & $0.2449$ & $0.1134$ & $0.2988$ & $0.1700$ \\
 & CORAL & 2000 & $0.0408$ & $0.1020$ & $0.0408$ & $0.0619$ & $0.0614$ & $0.0289$ \\
 & FreeMatch & 1500 & $0.1735$ & $0.6327$ & $0.2551$ & $0.4227$ & $0.3710$ & $0.2030$ \\
\midrule
purity at 50\% completeness & DANN & 4000 & $0.2979$ & $0.2382$ & $0.1984$ & $0.1889$ & $0.2308$ & $0.0495$ \\
 & BNM & 4000 & $0.0915$ & $0.2707$ & $0.0571$ & $0.0179$ & $0.1093$ & $0.1117$ \\
 & CORAL & 2000 & $0.0194$ & $0.0147$ & $0.0072$ & $0.0072$ & $0.0121$ & $0.0060$ \\
 & FreeMatch & 1500 & $0.0991$ & $0.3726$ & $0.0803$ & $0.1776$ & $0.1824$ & $0.1336$ \\
\bottomrule
\end{tabular}
\end{table}


% GENERATED FILE -- do not edit by hand.
% Rebuild:  python make_all.py (generator: make_dilute10_tables.py)
% Appendix K: the 10 t-intervals. runs/dilute10_analysis/results.json G and D.
\begin{table}[htbp]
\centering
\caption{\textbf{The ten paired $t$-intervals of the 1:10 experiment} at four decimals, folds 1--4, cleaned pool (df${}=3$). $G$ is diluted $-$ $\lambda=0$; its primary form reads the control at 4000 and decides the branch, and its matched-step form, a pre-declared sensitivity readout, reads the control at the arm's own step. $D$ is balanced $-$ diluted at the same read-off; it carries no verdict and includes the noise of a changed batch order. No multiplicity correction. $^{*}$Interval excludes zero.}
\label{tab:dl-contrasts4}
\footnotesize
\setlength{\tabcolsep}{4pt}
\begin{tabular}{llrlcl}
\toprule
Contrast & Steps & $\Delta$ TPR@$10^{-3}$ & 95\% CI & Signs & Role \\
\midrule
\multicolumn{6}{l}{\textit{$G$: diluted $-$ $\lambda=0$}} \\
DANN $-$ $\lambda=0$ & 4000 / 4000 & $+0.3350$$^{*}$ & $[+0.2913, +0.3787]$ & 4/4 & primary \\
BNM $-$ $\lambda=0$ & 4000 / 4000 & $+0.1941$$^{*}$ & $[+0.0029, +0.3854]$ & 4/4 & primary \\
CORAL $-$ $\lambda=0$ & 2000 / 4000 & $+0.0102$ & $[-0.0251, +0.0454]$ & 3/4 & primary \\
CORAL $-$ $\lambda=0$ & 2000 / 2000 & $-0.0206$ & $[-0.0740, +0.0329]$ & 3/4 & sensitivity \\
FreeMatch $-$ $\lambda=0$ & 1500 / 4000 & $+0.2533$ & $[-0.0042, +0.5107]$ & 4/4 & primary \\
FreeMatch $-$ $\lambda=0$ & 1500 / 1500 & $+0.2431$ & $[-0.0229, +0.5090]$ & 4/4 & sensitivity \\
\midrule
\multicolumn{6}{l}{\textit{$D$: balanced $-$ diluted}} \\
DANN & 4000 / 4000 & $-0.0206$ & $[-0.2840, +0.2427]$ & 2/4 & secondary \\
BNM & 4000 / 4000 & $+0.0257$ & $[-0.1076, +0.1591]$ & 2/4 & secondary \\
CORAL & 2000 / 2000 & $+0.3273$$^{*}$ & $[+0.1925, +0.4621]$ & 4/4 & secondary \\
FreeMatch & 1500 / 1500 & $+0.0793$ & $[-0.0800, +0.2385]$ & 3/4 & secondary \\
\bottomrule
\end{tabular}
\end{table}


\paragraph{The frozen stages.} Table~\ref{tab:dl-frozen} gives the frozen stages that the
unfrozen CORAL and FreeMatch arms start from, on the diluted target and in the main design, with
the \dlNbootWord{} bootstrap intervals, and splits FreeMatch's contrast against the control into
the parts gained before and after unfreezing. Table~\ref{tab:dl-heads} gives the norms of those
stages' final heads, read from the stored weights. CORAL's diluted frozen head has a
layer-normalisation gain about \dlCoralGainRatioWord{} times that of its counterpart in the
main design,
and FreeMatch's has not changed in the same way; we did not establish why.

% GENERATED FILE -- do not edit by hand.
% Rebuild:  python make_all.py (generator: make_dilute10_tables.py)
% Appendix K: frozen stages vs round 1 and the FreeMatch decomposition. runs/dilute10_analysis/results.json descriptive.
\begin{table}[htbp]
\centering
\caption{\textbf{The frozen stages on the 1:10 target against the balanced target}, and where FreeMatch's contrast is gained. Top: TPR at FPR${}=10^{-3}$ of each frozen stage's step-8000 head per fold (1:1 is the main design), cleaned pool; the difference is diluted $-$ balanced under the class-conditional paired bootstrap ($B=1000$), as in the frozen comparison. Bottom: FreeMatch's contrast against the control ($\lambda=0$ at 4000) split into the part reached by its frozen head (start $-$ $\lambda=0$) and the part gained after unfreezing (unfrozen $-$ start). AUROC: the mean over the four folds of the same heads' AUROC on the same pool.}
\label{tab:dl-frozen}
\footnotesize
\setlength{\tabcolsep}{1.5pt}
\begin{tabular}{llrrrrrrrl}
\toprule
Stage & Target & fold 1 & fold 2 & fold 3 & fold 4 & mean & AUROC & 1:10 $-$ 1:1 & 95\% CI \\
\midrule
FreeMatch & 1:1 (main design) & $0.4184$ & $0.4184$ & $0.3776$ & $0.4330$ & $0.4118$ & $0.9586$ & & \\
 & 1:10 & $0.2551$ & $0.3367$ & $0.2653$ & $0.3505$ & $0.3019$ & $0.9475$ & $-0.1099$ & $[-0.1634, -0.0689]$ \\
CORAL & 1:1 (main design) & $0.1837$ & $0.2245$ & $0.1837$ & $0.2474$ & $0.2098$ & $0.9065$ & & \\
 & 1:10 & $0.1939$ & $0.1939$ & $0.1735$ & $0.2165$ & $0.1944$ & $0.9001$ & $-0.0154$ & $[-0.0461, +0.0076]$ \\
\midrule
 & Target & start & frozen & unfrozen & total & & & & \\
\midrule
FreeMatch & 1:1 (main design) & $0.4118$ & $+0.2276$ & $+0.1049$ & $+0.3325$ & & & & \\
FreeMatch & 1:10 & $0.3019$ & $+0.1177$ & $+0.1355$ & $+0.2533$ & & & & \\
\bottomrule
\end{tabular}
\end{table}


% GENERATED FILE -- do not edit by hand.
% Rebuild:  python make_all.py (generator: make_dilute10_tables.py)
% Appendix K: frozen-stage head norms. runs/dilute10_head_norms/head_norms.json.
\begin{table}[htbp]
\centering
\caption{\textbf{The frozen-stage heads} that initialise the unfrozen CORAL and FreeMatch arms, read from their stored step-8000 state dicts: the $L_2$ norm of the layer-normalisation gain and of its bias, and $\|\mathrm{head}\|_2$ over every head parameter. A gain of one in every dimension has norm $\sqrt{640}=\dlUnitGain{}$. Each $\|\mathrm{head}\|_2$ equals the value its run recorded when it was used.}
\label{tab:dl-heads}
\footnotesize
\setlength{\tabcolsep}{4pt}
\begin{tabular}{lllrrrr}
\toprule
Head & Target & Norm & fold 1 & fold 2 & fold 3 & fold 4 \\
\midrule
CORAL & main design (1:1) & gain & $22.4$ & $33.1$ & $24.0$ & $23.4$ \\
 &  & bias & $63.6$ & $64.5$ & $63.7$ & $61.6$ \\
 &  & $\|\mathrm{head}\|_2$ & $67.5$ & $72.6$ & $68.2$ & $66.0$ \\
 & 1:10 & gain & $97.0$ & $103.7$ & $98.2$ & $102.2$ \\
 &  & bias & $92.6$ & $92.6$ & $90.2$ & $91.3$ \\
 &  & $\|\mathrm{head}\|_2$ & $134.2$ & $139.0$ & $133.4$ & $137.1$ \\
\midrule
FreeMatch & main design (1:1) & gain & $28.9$ & $28.9$ & $29.1$ & $28.8$ \\
 &  & bias & $0.7$ & $0.7$ & $0.8$ & $0.7$ \\
 &  & $\|\mathrm{head}\|_2$ & $29.3$ & $29.3$ & $29.5$ & $29.2$ \\
 & 1:10 & gain & $27.9$ & $27.9$ & $28.1$ & $27.9$ \\
 &  & bias & $0.9$ & $0.9$ & $0.8$ & $0.8$ \\
 &  & $\|\mathrm{head}\|_2$ & $28.2$ & $28.3$ & $28.4$ & $28.2$ \\
\bottomrule
\end{tabular}
\end{table}


\paragraph{The predictions.} Table~\ref{tab:dl-predictions} scores every component of our
\dlNPredWord{} predictions as written. Components that name a mechanism are not testable with
any stored quantity and are marked so.

% GENERATED FILE -- do not edit by hand.
% Rebuild:  python make_all.py (generator: make_dilute10_tables.py)
% Appendix K: prediction scoring. runs/dilute10_analysis/results.json predictions.
\begin{table}[htbp]
\centering
\caption{\textbf{Our four predictions for the 1:10 experiment, scored as written.} Each prediction, recorded before any run, is split into its components; where its wording admits two literal readings, both are scored. \emph{Rests on} gives the stored value each score reads; relative loss is $D$ / (balanced $-$ $\lambda=0$). No prediction was amended. Component 3d(ii) takes a maximum over \dlBnmPpN{} single-batch values and would be met by batch noise alone; we do not count it as support (Section~\ref{sec:results-dilute}).}
\label{tab:dl-predictions}
\scriptsize
\setlength{\tabcolsep}{4pt}
\begin{tabular}{l>{\TRAG}p{5.4cm}l>{\TRAG}p{5.0cm}}
\toprule
ID & Component & Score & Rests on \\
\midrule
1a & DANN degrades (D point $>$ 0) & failed & D = $-$0.0206 \\
1b(i) & G point $<$ +0.31 (as stated) & failed & G = +0.3350 \\
1b(ii) & G point $<$ reproduced balanced $-$ $\lambda=0$ & failed & G = +0.3350 vs +0.3144 \\
1c(i) & D is positive: point & failed & D = $-$0.0206 \\
1c(ii) & D is positive: interval excludes zero above & failed & D CI [$-$0.2840, +0.2427] \\
1d & G's interval (probably) includes zero & failed & G CI [+0.2913, +0.3787] \\
1e & branch COLLAPSES or SHRINKS & failed & branch SURVIVES \\
1e' & 'with a wide D' & not testable & 'wide' is not operationalised in part 1 \\
1f & mechanism (marginal alignment pulls lenses toward galaxies; prior collapses) & not testable & no stored quantity measures it \\
\midrule
2a & CORAL degrades (D point $>$ 0) & held & D = +0.3273 \\
2b(i) & less than DANN, absolute & failed & D CORAL +0.3273 vs DANN $-$0.0206 \\
2b(ii) & less than DANN, relative & failed & relative loss CORAL 0.9699 vs DANN $-$0.0657 \\
2c & G's interval excludes zero & failed & G CI [$-$0.0251, +0.0454] \\
2d & branch SHRINKS or SURVIVES & failed & branch COLLAPSES \\
2e & mechanism (moment matching a weaker pull) & not testable & no stored quantity \\
\midrule
3a & BNM degrades (D point $>$ 0) & held & D = +0.0257 \\
3b & the most of the four in relative terms & failed & relative loss DANN $-$0.0657, BNM 0.1171, CORAL 0.9699, FreeMatch 0.2384 \\
3c & G's interval may include zero & not testable & permissive wording; G CI [+0.0029, +0.3854] \\
3d(i) & predicted-positive fraction $>$ 0.09 at step 4000 & mixed 1/4 & step-4000 values f1 0.1055, f2 0.0586, f3 0.0742, f4 0.0352 \\
3d(ii) & predicted-positive fraction $>$ 0.09 at any logged step & held & max per fold f1 0.1055, f2 0.0938, f3 0.1016, f4 0.0938 \\
3e & mechanism (high-scoring negatives) & not testable & needs a new measurement \\
\midrule
4a & FreeMatch degrades (D point $>$ 0) & held & D = +0.0793 \\
4b(i) & less than BNM, absolute & failed & D FreeMatch +0.0793 vs BNM +0.0257 \\
4b(ii) & less than BNM, relative & failed & relative loss FreeMatch 0.2384 vs BNM 0.1171 \\
4c & G's interval still excludes zero & failed & G CI [$-$0.0042, +0.5107] \\
4d(i) & head start shrinks: point & held & diluted frozen 0.3019 vs main-design frozen 0.4118 \\
4d(ii) & head start shrinks: bootstrap interval wholly below zero & held & diluted $-$ main design = $-$0.1099 [$-$0.1634, $-$0.0689] \\
4e & mechanism (uniform fairness target; adaptive thresholds) & not testable & no stored quantity \\
\bottomrule
\end{tabular}
\end{table}


\paragraph{Other test pools.} Table~\ref{tab:dl-sensitivity} repeats the \dlNtWord{} contrasts on
the original test pool and on the reduced pool. On the reduced pool every contrast keeps its
zero-inclusion. On the original pool \dlOrigMatchWord{} of the \dlNtWord{} do; both forms of
FreeMatch's $G$ exclude zero there, the primary at $\dlGFmOrig{}$
$[\dlGFmOrigLo{}, \dlGFmOrigHi{}]$, which Section~\ref{sec:results-dilute} notes beside the
near-boundary pair.

% GENERATED FILE -- do not edit by hand.
% Rebuild:  python make_all.py (generator: make_dilute10_tables.py)
% Appendix K: pool sensitivities. runs/dilute10_analysis/results.json sensitivities (orig, D2).
\begin{table}[htbp]
\centering
\caption{\textbf{The ten contrasts of the 1:10 experiment on the original and the reduced test pools}, beside whether the cleaned pool's interval excludes zero. The original pool keeps every drawn test negative (Section~\ref{sec:testpool}); the reduced pool also removes the test lenses that appear in training or validation. Match: the same zero-inclusion as on the cleaned pool. Paired $t$-intervals (df${}=3$).}
\label{tab:dl-sensitivity}
\footnotesize
\setlength{\tabcolsep}{2.5pt}
\begin{tabular}{llcrcrc}
\toprule
Contrast & Steps & Cleaned & Original pool & Match & Reduced pool & Match \\
\midrule
$G$ DANN & 4000 / 4000 & excl. & $+0.2814$ $[+0.2388, +0.3240]$ & yes & $+0.3336$ $[+0.2846, +0.3826]$ & yes \\
$G$ BNM & 4000 / 4000 & excl. & $+0.1610$ $[+0.0471, +0.2749]$ & yes & $+0.2015$ $[+0.0036, +0.3994]$ & yes \\
$G$ CORAL & 2000 / 4000 & incl. & $+0.0333$ $[-0.0006, +0.0672]$ & yes & $-0.0009$ $[-0.0601, +0.0584]$ & yes \\
$G$ CORAL & 2000 / 2000 & incl. & $-0.0052$ $[-0.0561, +0.0457]$ & yes & $-0.0248$ $[-0.0901, +0.0405]$ & yes \\
$G$ FreeMatch & 1500 / 4000 & incl. & $+0.2099$ $[+0.0051, +0.4147]$ & \textbf{no} & $+0.2457$ $[-0.0035, +0.4948]$ & yes \\
$G$ FreeMatch & 1500 / 1500 & incl. & $+0.1997$ $[+0.0071, +0.3923]$ & \textbf{no} & $+0.2491$ $[-0.0187, +0.5169]$ & yes \\
$D$ DANN & 4000 / 4000 & incl. & $-0.0002$ $[-0.2233, +0.2230]$ & yes & $-0.0191$ $[-0.2788, +0.2406]$ & yes \\
$D$ BNM & 4000 / 4000 & incl. & $+0.0310$ $[-0.0915, +0.1534]$ & yes & $+0.0220$ $[-0.1083, +0.1524]$ & yes \\
$D$ CORAL & 2000 / 2000 & excl. & $+0.2660$ $[+0.1065, +0.4255]$ & yes & $+0.3364$ $[+0.2066, +0.4662]$ & yes \\
$D$ FreeMatch & 1500 / 1500 & incl. & $+0.0715$ $[-0.0542, +0.1972]$ & yes & $+0.0805$ $[-0.0782, +0.2392]$ & yes \\
\bottomrule
\end{tabular}
\end{table}


\paragraph{Run bookkeeping.} Table~\ref{tab:dl-bookkeeping} gives each half-precision run's
skipped and applied optimiser updates, the realised lens fraction of every run's target batches
against the pool's expected value, and the firings of CORAL's advisory feature-norm guard. The
no-verdict rule of the third part of the pre-registration fires below half the budget's updates;
the lowest applied fraction is \dlCoralAppliedPct{}\% for CORAL and \dlFmAppliedPct{}\% for
FreeMatch. CORAL's skips are its first \dlCoralSkipMin{}--\dlCoralSkipMax{} steps on every fold,
after which none is skipped; FreeMatch skips \dlFmSkipMin{}--\dlFmSkipMax{} per fold. The lens
fractions of DANN and BNM sit below those of the other arms because their runs keep the
half-lens holdout of the balanced design (Section~\ref{sec:labels}).

% GENERATED FILE -- do not edit by hand.
% Rebuild:  python make_all.py (generator: make_dilute10_tables.py)
% Appendix K: run bookkeeping. runs/dilute10_analysis/results.json descriptive.bookkeeping; runs/dilute10/coral_nb3/fold*/dilute10_feature_norm_advisory.json.
\begin{table}[htbp]
\centering
\caption{\textbf{Run bookkeeping of the 1:10 experiment.} GradScaler skips and the fraction of the budget's optimiser updates applied (half-precision runs; the CORAL frozen stage ran in full precision on CPU, and the BNM runs and the FreeMatch frozen runs did not record skips); the lens fraction of the target batches, its run mean $\pm$ sd over the run's logged batches, against the pool's expected value; and CORAL's advisory feature-norm guard: the norm it reads (the mean $L_2$ norm of a batch's feature vectors, source and target) at step 1 and its smallest value over the run, source or target, against the guard's threshold, and its firings.}
\label{tab:dl-bookkeeping}
\footnotesize
\setlength{\tabcolsep}{2.5pt}
\begin{tabular}{llrrrr}
\toprule
Run & Quantity & fold 1 & fold 2 & fold 3 & fold 4 \\
\midrule
DANN & skips / applied & 3 / $0.9992$ & 6 / $0.9985$ & 3 / $0.9992$ & 3 / $0.9992$ \\
CORAL & skips / applied & 9 / $0.9955$ & 10 / $0.9950$ & 9 / $0.9955$ & 9 / $0.9955$ \\
FreeMatch & skips / applied & 2 / $0.9987$ & 2 / $0.9987$ & 3 / $0.9980$ & 3 / $0.9980$ \\
\midrule
DANN & lens fraction, mean $\pm$ sd & $0.0834 \pm 0.0164$ & $0.0833 \pm 0.0169$ & $0.0833 \pm 0.0167$ & $0.0834 \pm 0.0163$ \\
 & expected & $0.0834$ & $0.0834$ & $0.0834$ & $0.0834$ \\
BNM & lens fraction, mean $\pm$ sd & $0.0834 \pm 0.0164$ & $0.0833 \pm 0.0169$ & $0.0833 \pm 0.0167$ & $0.0834 \pm 0.0163$ \\
 & expected & $0.0834$ & $0.0834$ & $0.0834$ & $0.0834$ \\
CORAL & lens fraction, mean $\pm$ sd & $0.0909 \pm 0.0177$ & $0.0909 \pm 0.0173$ & $0.0909 \pm 0.0181$ & $0.0909 \pm 0.0176$ \\
 & expected & $0.0909$ & $0.0909$ & $0.0909$ & $0.0909$ \\
FreeMatch & lens fraction, mean $\pm$ sd & $0.0909 \pm 0.0177$ & $0.0909 \pm 0.0172$ & $0.0909 \pm 0.0181$ & $0.0909 \pm 0.0176$ \\
 & expected & $0.0909$ & $0.0909$ & $0.0909$ & $0.0909$ \\
CORAL frozen & lens fraction, mean $\pm$ sd & $0.0909 \pm 0.0177$ & $0.0909 \pm 0.0176$ & $0.0909 \pm 0.0177$ & $0.0909 \pm 0.0173$ \\
 & expected & $0.0909$ & $0.0909$ & $0.0909$ & $0.0909$ \\
FreeMatch frozen & lens fraction, mean $\pm$ sd & $0.0909 \pm 0.0177$ & $0.0909 \pm 0.0176$ & $0.0909 \pm 0.0176$ & $0.0909 \pm 0.0173$ \\
 & expected & $0.0909$ & $0.0909$ & $0.0909$ & $0.0909$ \\
\midrule
CORAL & norm at step 1, source & $129.8$ & $133.1$ & $126.3$ & $131.8$ \\
 & norm at step 1, target & $149.1$ & $158.6$ & $150.1$ & $154.7$ \\
 & smallest norm & $121.0$ & $121.5$ & $119.8$ & $117.9$ \\
 & guard threshold & $110$ & $110$ & $110$ & $110$ \\
 & advisory firings & 2000 & 2000 & 2000 & 2000 \\
\bottomrule
\end{tabular}
\end{table}


\paragraph{The memorisation diagnostic.} Table~\ref{tab:dl-memorisation} gives the DANN
$-$ $\lambda=0$ analysis of Section~\ref{sec:negative} per fold. It reads the stored
checkpoints of the main design's runs, not the diluted ones.

% GENERATED FILE -- do not edit by hand.
% Rebuild:  python make_all.py (generator: make_dilute10_tables.py)
% Appendix K: DANN - lambda=0 memorisation diagnostic. runs/dann_lambda0_holdout_gap/results.json.
\begin{table}[htbp]
\centering
\caption{\textbf{The memorisation diagnostic for DANN and its control}, a retrospective analysis declared before it was computed. Each value is the gap in the normalised nuclear norm of the predictions between the target slice the arm trains on and the held-out slice, \mzHoldout{} images each, computed without augmentation on the stored checkpoints: at the read-off step (primary) and averaged over the \mzNckpt{} checkpoints (secondary). The maxima over checkpoints are descriptive only, because a maximum of noisy reads is positive even when no gap exists. These values are not comparable with BNM's logged gap, which used one augmented view per image.}
\label{tab:dl-memorisation}
\footnotesize
\setlength{\tabcolsep}{4pt}
\begin{tabular}{llrrrrrl}
\toprule
Reading & Arm & fold 1 & fold 2 & fold 3 & fold 4 & mean & 95\% CI \\
\midrule
step 4000 & DANN & $+0.0237$ & $+0.0250$ & $+0.1580$ & $+0.2288$ & $+0.1089$ & \\
 & $\lambda=0$ & $+0.2119$ & $-0.1226$ & $-0.0396$ & $-0.0089$ & $+0.0102$ & \\
 & DANN $-$ $\lambda=0$ & $-0.1882$ & $+0.1476$ & $+0.1977$ & $+0.2377$ & $+0.0987$ & $[-0.2113, +0.4086]$ \\
\midrule
mean over steps & DANN & $+0.0803$ & $+0.0641$ & $+0.0874$ & $+0.1247$ & $+0.0891$ & \\
 & $\lambda=0$ & $+0.1050$ & $-0.0671$ & $+0.0748$ & $-0.0676$ & $+0.0113$ & \\
 & DANN $-$ $\lambda=0$ & $-0.0248$ & $+0.1313$ & $+0.0126$ & $+0.1924$ & $+0.0779$ & $[-0.0832, +0.2390]$ \\
\midrule
maximum & DANN & $+0.1932$ & $+0.1732$ & $+0.1841$ & $+0.2739$ & & \\
 & \quad at step & 2500 & 3000 & 750 & 2250 & & \\
 & $\lambda=0$ & $+0.2894$ & $+0.0698$ & $+0.1994$ & $+0.0943$ & & \\
 & \quad at step & 3000 & 1000 & 500 & 2000 & & \\
\bottomrule
\end{tabular}
\end{table}

